\chapter{Validazione e risultati ottenuti}
In questo capitolo il sistema viene validato su dati sperimentali reali, verificando come le funzionalità implementate rispondano ai requisiti individuati nel \cref{cap:software_esistente_requisiti}.
Dopo una descrizione dei dati utilizzati, vengono presentate le principali funzionalità dell'event display attraverso schermate commentate, per poi concludere con una verifica dei requisiti e una discussione dei risultati ottenuti.

Gli esempi mostrati nelle sezioni successive fanno riferimento alla run \texttt{5262}, evento \texttt{112104442}, acquisita il 16 novembre 2022 durante la Run 3 di LHC.
I dati di questa run appartengono al dataset utilizzato dalla collaborazione per la prima osservazione di interazioni di neutrini senza muoni nello stato finale, pubblicata su \textit{Physical Review Letters} nel giugno 2025 \cite{SND_Observation}.
La scelta di questo evento consente inoltre un confronto diretto con il software preesistente, del quale era già stato mostrato l'output in \cref{fig:vecchio_display}.

% ====================================================
%
% 6.1 - FUNZIONALITA' DEL SISTEMA
%
% ====================================================
\section{Funzionalità del sistema}
Vengono ora illustrate le principali funzionalità dell'event display attraverso esempi su dati reali.

% ====================================================
% 6.1.1 - SELEZIONE DELL'EVENTO
% ====================================================
\subsection{Selezione dell'evento}
La selezione dell'evento da visualizzare avviene attraverso due schermate successive.
Nella prima l'utente inserisce il numero di run (\cref{fig:scelta_run}); nella seconda vengono mostrati i metadati della run caricata e viene richiesto il numero di evento, insieme alle configurazioni per il clustering e la colorazione delle hit (\cref{fig:scelta_evento}), che verranno illustrate successivamente.

Le stesse schermate sono inoltre accessibili durante l'interazione tramite apposite voci di menu, consentendo di cambiare i dati visualizzati senza riavviare il software.
Per rendere la navigazione di eventi sequenziali il più agevole possibile, durante l'interazione è possibile scorrere gli eventi consecutivi in avanti e indietro rispettivamente utilizzando i pulsanti \texttt{K} e \texttt{J} della tastiera.

\begin{figure}[H]
    \centering
    \begin{subfigure}[b]{0.49\textwidth}
        \centering
        \includegraphics[width=\textwidth]{cap6/SceltaRun}
        \caption{Scelta run}
        \label{fig:scelta_run}
    \end{subfigure}
    \hfill
    \begin{subfigure}[b]{0.49\textwidth}
        \centering
        \includegraphics[width=\textwidth]{cap6/SceltaEvento}
        \caption{Scelta evento}
        \label{fig:scelta_evento}
    \end{subfigure}
    \caption{Selezione dei dati da visualizzare}
\end{figure}

Qualora il file di geometria predefinito non venga trovato, il sistema notifica l'errore indicando esplicitamente il file richiesto per la run selezionata (\cref{fig:errore_geometria}).
Consente poi di selezionarne manualmente uno alternativo tramite un file browser (\cref{fig:scelta_geometria}), riprendendo l'inizializzazione senza richiedere il riavvio del software.
Il file browser rimane inoltre accessibile durante l'interazione per caricare geometrie alternative.

\begin{figure}[ht]
    \centering
    \begin{subfigure}[b]{0.49\textwidth}
        \centering
        \includegraphics[width=\textwidth]{cap6/GeometriaRichiesta}
        \caption{Geometria richiesta}
        \label{fig:errore_geometria}
    \end{subfigure}
    \hfill
    \begin{subfigure}[b]{0.49\textwidth}
        \centering
        \includegraphics[width=\textwidth]{cap6/GeometriaManuale}
        \caption{Scelta manuale geometria}
        \label{fig:scelta_geometria}
    \end{subfigure}
    \caption[Geometria predefinita non trovata]{Gestione del file di geometria non trovato nel percorso predefinito}
\end{figure}

Questa modalità di selezione soddisfa il requisito funzionale di selezionare run ed evento tramite interfaccia, eliminando la necessità di conoscere il sistema di archiviazione EOS e i comandi da shell illustrati nel \cref{lst:vecchio_display}.

% ====================================================
% 6.1.2 - CONFIGURAZIONE CLUSTER E COLORAZIONE
% ====================================================
\subsection{Configurazione cluster e colorazione}
Le opzioni di clustering e colorazione (\cref{fig:cluster_config}) sono accessibili nella schermata di selezione dell'evento e possono essere modificate in qualsiasi momento durante l'interazione, anche mantenendo invariato il numero di evento.

\begin{figure}[ht]
    \centering
    \includegraphics[width=\textwidth]{cap6/ConfigurazioniCluster}
    \caption[Configurazione dei cluster]{Configurazione dell'algoritmo di clustering e della scala di colore nella schermata di selezione evento}
    \label{fig:cluster_config}
\end{figure}

Per il clustering, i parametri configurabili sono il numero minimo di hit necessarie a formare un cluster e la distanza massima consentita tra hit adiacenti.
Raggruppare le hit vicine riduce l'affollamento visivo della scena, preservando le informazioni fisicamente rilevanti.
Data la differente risoluzione di ogni sottorivelatore, studiata nel \cref{cap:snd}, è messo a disposizione un set di parametri per ognuno di essi.
Durante l'interazione i parametri utilizzati sono riportati in una finestra, visualizzata in \cref{fig:visualizza_cluster}.
Per evitare di doverli impostare singolarmente ad ogni utilizzo del display, sono presenti dei valori predefiniti, scelti dalla collaborazione durante la fase di analisi.

Per la colorazione, è possibile scegliere la grandezza fisica di riferimento -- energia o tempo -- e la scala cromatica da applicare: lineare, logaritmica o monocromatica.
Questa flessibilità consente di evidenziare aspetti diversi dell'evento a seconda dell'analisi in corso.
È possibile abilitare la visualizzazione della scala (mostrata in \cref{fig:scala_colore}), così da semplificare l'associazione di ogni colore al valore di riferimento.

\begin{figure}[htb]
    \centering
    \hfill
    \begin{subfigure}[b]{0.3\textwidth}
        \centering
        \includegraphics[width=\textwidth]{cap6/ConfigurazioneCluster}
        \caption{Configurazioni del clustering}
        \label{fig:visualizza_cluster}
    \end{subfigure}
    \hfill
    \begin{subfigure}[b]{0.35\textwidth}
        \centering
        \includegraphics[width=\textwidth]{cap6/ScalaColore}
        \caption{Scala colore}
        \label{fig:scala_colore}
    \end{subfigure}
    \hfill
    \caption{Visualizzazione configurazioni dei dati caricati}
\end{figure}

Nel complesso, le due opzioni soddisfano rispettivamente i requisiti di selezionare le possibili configurazioni dei dati da caricare e di visualizzare l'energia depositata in modo coerente con i dati letti tramite ROOT.

% ====================================================
% 6.1.3 - INTERAZIONE CON LA SCENA
% ====================================================
\subsection{Interazione con la scena}
Una volta caricati i dati, è possibile interagire liberamente con la scena tridimensionale tramite rotazione, traslazione e zoom, soddisfacendo il requisito di osservare la geometria del rivelatore da prospettive arbitrarie.
I controlli sono accessibili tramite il puntatore e lo scroll del mouse, con il tasto Shift come modificatore per la traslazione.
Sono inoltre disponibili quattro viste predefinite -- frontale, laterale, dall'alto e isometrica -- richiamabili tramite appositi tasti della tastiera, ciascuna delle quali attiva automaticamente la proiezione ortografica.
In \cref{fig:viste_predefinite} sono riportate le viste dall'alto e laterale, corrispondenti alle proiezioni già presenti nel software preesistente.

\begin{figure}[htb]
    \centering
    \begin{subfigure}[b]{0.8\textwidth}
        \centering
        \includegraphics[width=\textwidth]{cap6/VistaAlto}
        \caption{Vista dall'alto}
    \end{subfigure}
    \hfill
    \begin{subfigure}[b]{0.8\textwidth}
        \centering
        \includegraphics[width=\textwidth]{cap6/VistaLaterale}
        \caption{Vista laterale}
    \end{subfigure}
    \caption{Attivazione delle viste predefinite}
    \label{fig:viste_predefinite}
\end{figure}

\subsubsection{Visibilità dei nodi}
Il pannello \texttt{Scene} (\cref{fig:pannello_scene}), accessibile in qualsiasi momento durante l'interazione, espone l'albero dei nodi della scena diviso in due sezioni: la geometria del rivelatore (\texttt{SND}) e le hit dell'evento (\texttt{HITS}).
Per ciascun nodo è disponibile una checkbox che ne abilita o disabilita la visualizzazione a runtime, consentendo di isolare singoli sottorivelatori senza ricaricare i dati.

\begin{figure}[htb]
    \centering
    \hfill
    \begin{subfigure}[b]{0.45\textwidth}
        \centering
        \includegraphics[width=\textwidth]{cap6/NodiGeometria}
        \caption{Albero della geometria}
    \end{subfigure}
    \hfill
    \begin{subfigure}[b]{0.45\textwidth}
        \centering
        \includegraphics[width=\textwidth]{cap6/NodiHit}
        \caption{Albero delle hit}
    \end{subfigure}
    \hfill
    \caption[Pannello degli alberi dei nodi]{Le due sezioni del pannello \texttt{Scene}: la geometria del rivelatore e le hit dell'evento}
    \label{fig:pannello_scene}
\end{figure}

% ====================================================
% 6.1.4 - ILLUMINAZIONE E TRASPARENZA
% ====================================================
\subsection{Illuminazione e trasparenza}
Un menu dedicato consente all'utente di attivare e disattivare l'illuminazione e la trasparenza, oltre a configurarne i parametri: per la trasparenza è possibile regolare separatamente l'opacità dei bordi e delle facce dei poligoni; per l'illuminazione è possibile modificare le proprietà delle sorgenti luminose puntiformi.

Il confronto tra la scena illuminata e non è apprezzabile nella \cref{fig:confronto_illuminazione}: sebbene non prevista tra i requisiti iniziali, senza illuminazione le hit appaiono come solidi uniformemente colorati, mentre, con l'illuminazione attiva, il calcolo della luce incidente rende distinguibili le singole facce, migliorando la percezione della profondità e della forma degli oggetti nella scena.

\begin{figure}[H]
    \centering
    \begin{subfigure}[b]{0.49\textwidth}
        \centering
        \includegraphics[width=\textwidth]{cap6/NonIlluminato}
        \caption{Illuminazione disattivata}
    \end{subfigure}
    \hfill
    \begin{subfigure}[b]{0.49\textwidth}
        \centering
        \includegraphics[width=\textwidth]{cap6/Illuminato}
        \caption{Illuminazione attivata}
    \end{subfigure}
    \caption[Confronto tra illuminazione attivata e disattivata]{Confronto di parte del rivelatore con illuminazione attivata e disattivata}
    \label{fig:confronto_illuminazione}
\end{figure}

% ====================================================
%
% 6.2 - VERIFICA DEI REQUISITI
%
% ====================================================
\section{Verifica dei requisiti}
Sulla base delle funzionalità presentate nella sezione precedente, tutti i requisiti funzionali individuati nel \cref{cap:software_esistente_requisiti} risultano soddisfatti.
La selezione di run ed evento avviene tramite schermate di inserimento guidato, eliminando la necessità di conoscere il sistema di archiviazione EOS e i comandi da shell del software preesistente.
L'interazione con la scena è gestita tramite mouse e tastiera, con viste predefinite richiamabili istantaneamente.
La visualizzazione dell'energia depositata è supportata attraverso la colorazione delle hit per energia o tempo, con scala cromatica configurabile.
Infine, i parametri di clustering e la modalità di colorazione costituiscono le configurazioni dei dati selezionabili dall'utente.

Anche i requisiti non funzionali sono rispettati.
La fluidità dell'interazione è garantita dalle ottimizzazioni discusse nella \cref{subsec:ottimizzazioni}: il precalcolo delle matrici di modellazione e la riduzione delle chiamate al driver grafico consentono di raggiungere 40-60 FPS con illuminazione e trasparenza attive, un valore in linea con lo standard di fluidità delle applicazioni interattive moderne.
Il ridimensionamento della finestra è gestito tramite le callback di GLFW, che aggiornano automaticamente l'area di disegno e il rapporto d'aspetto della proiezione, rendendo l'interfaccia adattabile a qualsiasi configurazione dello schermo.
La navigazione tra eventi consecutivi è resa agevole dai tasti \texttt{J} e \texttt{K}, che consentono di scorrere gli eventi in avanti e indietro senza dover reinserire i dati ogni volta.
L'interazione intuitiva è infine garantita dalla barra menu, che espone tutte le funzionalità disponibili riportando accanto a ciascuna la scorciatoia da tastiera corrispondente, così da non richiedere all'utente la memorizzazione di comandi.

Nella \cref{fig:schermata_completa} è riportata l'intera schermata del software mostrata all'utente durante l'interazione con la scena, in cui sono visibili la barra menu con un sottomenu aperto e le indicazioni dello stato corrente e il numero di FPS, i pannelli informativi con i metadati dell'evento e la palette utilizzata, e la finestra contenente le impostazioni grafiche del software.

\begin{figure}[ht]
    \centering
    \includegraphics[width=\textwidth]{cap6/Interfaccia}
    \caption{Interfaccia completa dell'event display}
    \label{fig:schermata_completa}
\end{figure}

% ====================================================
% CONCLUSIONE CAPITOLO 6
% ====================================================
\bigskip
I risultati ottenuti verranno ripresi nelle conclusioni, dove saranno discussi i limiti ancora presenti e i possibili sviluppi futuri del sistema.